\documentclass[11pt]{article}
\usepackage[margin=.88in]{geometry}
\usepackage[T1]{fontenc}
\usepackage{lmodern,microtype,amsmath,amssymb,graphicx,booktabs,tabularx}
\usepackage{xcolor}
\usepackage[colorlinks=true,linkcolor=blue!40!black,urlcolor=blue!50!black]{hyperref}
\setlength{\parindent}{0pt}
\setlength{\parskip}{.4em}
\setlength{\emergencystretch}{1em}
\graphicspath{{images/}}
\title{Conservative cluster encounters with a finite-work stock contact}
\author{Research calculation for the Entropic Scalar EFT project}
\date{23 October 2026}
\begin{document}
\maketitle
\begin{abstract}
A three-dimensional encounter evolves kinetic stock, stellar tracers and
pressure-supported baryonic gas in one gravitational potential. A selected
local stock--baryon interaction changes relative radial velocities, supplies
the opposite impulse to the partner, and records work and heat in finite
internal registers. Its symmetric clock gate is inherited from the preceding
effective-coupling construction. Seventeen runs test zero, weak and strong
contact, depleted registers, isolated systems, a collisionless-baryon control,
particle resolution, timestep, sampling and interaction range. Weak contact
preserves two total-mass concentrations associated with the stellar populations
while the gas is displaced. Strong contact changes the encounter and also
restructures the isolated starting halo. The latter result requires an
equilibrium-selection calculation before the strong-contact model can be used
predictively. Mass, linear momentum, angular momentum and the complete
Newtonian energy ledger are checked. Eight additional Lean certificates verify
the local impulse, work and gas-energy identities. Initial profiles, contact
target, physical rate, range and register allocation remain specified effective
inputs. No observed cluster likelihood has been fitted.
\end{abstract}

\section{The cluster question}
The cluster reciprocity audit identified a conservative route in which
positive inertial stock, stars and gas source a common potential. Its remaining
dynamical requirement was an explicit redistribution interaction with reaction
forces, energy transfer and a timescale. The present calculation implements
one such interaction in a live encounter.

The observational motivation is the separation of gas and lensing mass in
colliding clusters\cite{Clowe}. A comparison with such data requires evolving
both collisionless trajectories and gas stresses. Even the relation between a
gas shock speed and a subcluster speed depends on the encounter dynamics
\cite{Springel}. The calculation below uses declared idealized initial profiles;
an observed final mass map supplies none of its initial particle positions,
velocities or component masses.

The phase label of baryonic gas does not change its gravitational charge.
All surviving stock is counted once. Changes in projected mass arise from
transport, and the same total source determines the Newtonian dynamics and
the projected mass that would enter a weak-field lensing calculation.

\section{Common gravity and baryonic gas}
For particles of positive fixed mass $m_i$, use
\begin{equation}
 U_G=-G\sum_{i<j}\frac{m_i m_j}{\sqrt{|\mathbf x_i-\mathbf x_j|^2+\epsilon_G^2}},
 \qquad
 \mathbf a_i=-G\sum_{j\ne i}
 \frac{m_j(\mathbf x_i-\mathbf x_j)}{(|\mathbf x_i-\mathbf x_j|^2+\epsilon_G^2)^{3/2}}.
\end{equation}
This identical softened interaction acts on stock, stars and gas. Stock and
stellar elements retain separate trajectories and can multistream.
Cosmological exchange is neglected over this encounter, as in the cluster
audit's local dynamical limit.

Gas obeys an ideal $\gamma_g=5/3$ pressure law, with conservative smoothed
particle hydrodynamics (SPH)\cite{Cossins}. For fixed smoothing length $h_g$,
\begin{align}
 \rho_i&=\sum_{j\in g}m_jW_{h_g}(r_{ij}),\qquad
 A_i=\frac{P_i}{\rho_i^2}=\frac{(\gamma_g-1)u_i}{\rho_i},\\
 \dot{\mathbf v}_i\big|_g
 &=-\sum_{j\in g}m_j(A_i+A_j+\Pi_{ij})\nabla_iW_{ij},\\
 \dot u_i&=\sum_{j\in g}m_j(A_i+\Pi_{ij}/2)
 (\mathbf v_i-\mathbf v_j)\cdot\nabla_iW_{ij}.
\end{align}
The cubic-spline kernel has support $2h_g$. The symmetric Monaghan viscosity
uses $\alpha=1$, $\beta=2$ and acts on approaching pairs. It converts resolved
kinetic energy into gas heat. The paired pressure and internal-energy terms
cancel exactly in the instantaneous total-energy derivative.

This gas is baryonic plasma. The $\gamma=4/3$ warm capacity fluid of the
preceding note is not substituted for the kinetic stock. Keeping the latter
kinetic permits overlapping streams. The full spatial clock and floor-fluid
stress are not evolved here; the source-clock connection enters through
effective internal work registers and their transition gate.

\section{A reciprocal extension to moving carriers}
The earlier radial contact used one fixed spherical carrier. During an
encounter, a local baryonic partner moves and must receive the reaction. Let
$i$ be a stock element and $j$ a baryonic element. Define
\begin{equation}
 \mathbf n=\frac{\mathbf x_i-\mathbf x_j}{r_{ij}},\qquad
 w=(\mathbf v_i-\mathbf v_j)\cdot\mathbf n,\qquad
 \mu_{ij}=\frac{m_i m_j}{m_i+m_j}.
\end{equation}
For a proposed new relative radial speed $w'$, the updates are
\begin{align}\label{eq:kicks}
 \mathbf v_i'&=\mathbf v_i+\frac{m_j}{m_i+m_j}(w'-w)\mathbf n,\\
 \mathbf v_j'&=\mathbf v_j-\frac{m_i}{m_i+m_j}(w'-w)\mathbf n.
\end{align}
The centre-of-mass velocity and relative tangential velocity are unchanged.
The impulse is central, so total linear and angular momenta are unchanged.
The exact mechanical work is
\begin{equation}\label{eq:kinetic}
 \Delta K_{ij}=\frac{\mu_{ij}}2(w'^2-w^2).
\end{equation}
The gravitational energy is unchanged during this instantaneous velocity
update. Locality is implemented with a finite kernel of support $2h_c$;
the range remains an effective input requiring a continuum or physical
matching calculation.

Each stock element proposes a contact at rate $\nu$. Available baryonic
partners are drawn with weights $m_jW_{h_c}(r_{ij})$, normalized over local
partners. An element with no partner has no event. Both gas and stars are
eligible, including elements from the other initial cluster. Initial cohort
labels have no effect on interaction eligibility.

\subsection{Target and clock gate}
The trial relative-velocity density retains the earlier two-exponential
shape, now applied to relative kinetic energy:
\begin{equation}
 h(q)=0.7e^{-q/0.25}+0.3e^{-q},\qquad
 q=\frac{|\mathbf v_i-\mathbf v_j|^2}{2e_*},\qquad e_*=0.2.
\end{equation}
At fixed relative tangential speed $w_t$, the conditional Gaussian-mixture
weights for $w'$ are proportional to
\begin{equation}
 b_a\sqrt{T_a}\exp[-w_t^2/(2e_*T_a)],
 \qquad (b_1,b_2)=(0.7,0.3),\quad(T_1,T_2)=(0.25,1).
\end{equation}
This is a selected multicarrier target. Its stationary distribution in a
self-gravitating medium has not been derived. The earlier theorem for an
orbital density $h(E,\ell)$ around one spherical carrier does not establish
stationarity of this pair-relative target. Likewise, the conserved angular
quantity in~\eqref{eq:kicks} is the pair's angular momentum; individual
orbital angular labels around several different centres need not be fixed.

For a finite event, set
\begin{equation}
 Q=\Theta_{ij}\log\frac{h(q')}{h(q)},\qquad
 \Theta_{ij}=10^{-6}\mu_{ij}e_*,\qquad
 W=\Delta K_{ij}+Q.
\end{equation}
Heat $Q$ enters the bath ledger; work $W$ leaves the stock element's internal
work register $U_i$. Thus
\begin{equation}\label{eq:ledger}
 U_i'=U_i-W,\qquad E_{B,i}'=E_{B,i}+Q,\qquad
 \Delta(K+U+E_B)=0.
\end{equation}
Each register has floor $K_i=m_i e_*$. The homogeneous clock-speed function
$v_c(y)$ from the preceding action, with $v_c(1)=0$ and $0\le v_c\le1$,
gives the symmetric acceptance factor
\begin{equation}\label{eq:gate}
 g_{ij}=v_c(U_i/K_i)v_c(U_i'/K_i).
\end{equation}
Proposals with $U_i'<K_i$ or $E_{B,i}'<0$ are unavailable. No negative
state is created and subsequently clipped.

The baseline internal allocation is $U_i=K_i(1+1.2)$ and
$E_{B,i}=10^{-4}K_i$. Separate controls use surplus $10^{-4}K_i$ or exactly
zero surplus. The small dimensionless $\Theta_{ij}$ is a thermostat choice;
it is not a conversion of the Gibbons--Hawking temperature to cluster units.
The finite bath ledger is kept nonnegative, while its temperature is held
fixed. A native bath spectrum and heat capacity remain unmodelled.

The work and bath registers are internal energies of the existing elements.
No additional rest-mass halo is introduced. Their contribution to gravitational
mass is neglected at Newtonian order. A covariant identification with the
capacity medium, including its spatial stress and inertial moments, remains
required. The present result certifies the classical event accounting under
this effective identification.

\section{Initial conditions and comparison protocol}
Units satisfy $G=M_1=a_1=1$. The second cluster has $M_2=0.3$ and
$a_2=0.3^{1/3}$. Each component begins with a Plummer density profile,
truncated at $6a_c$. Stock and stars use the isotropic Plummer velocity
distribution. Gas starts with zero internal bulk flow and specific energy
\begin{equation}
 u_g(r)=\frac{GM_c}{4\sqrt{r^2+a_c^2}},
\end{equation}
which gives hydrostatic support for a coextensive gas fraction in the
unsoftened total Plummer potential. Truncation, smoothing and finite sampling
make this an approximate equilibrium; isolated controls quantify the drift.
Antipodal phase sampling sets component centre-of-mass position, velocity
and initial spin exactly to their intended values.

The surviving stock-to-original-baryon ratio is the saved cohort convention
$k=5.1760752$. The mass fractions are
\begin{equation}
 f_C=0.8380849,\qquad f_*=0.02,\qquad f_g=0.1419151.
\end{equation}
They are allocations supplied to the encounter, not masses inferred from a
final lensing map. The initial centre separation is $(8,0.6,0)$ and the
relative velocity is $(-0.8,0,0)$. The centre of mass is at rest.

The baseline has 384 stock, 96 stellar and 192 gas elements per cluster,
or 1344 in total. Higher-resolution encounters double all three counts.
The standard gravity and gas smoothing lengths are $(\epsilon_G,h_g)=(0.1,0.45)$;
the finer calculation uses $(0.08,0.36)$. The contact scale stays $h_c=0.45$
during this refinement. One sensitivity run reduces it to $0.3$.
Time steps are $0.01$ and $0.005$. Implicit final-kick iteration treats
velocity-dependent gas viscosity, and contact events follow the mechanical
step. Distinct time steps have distinct random contact histories.

The 17 completed runs comprise 14 encounters and three isolated controls.
Encounter tests include $\nu=0$, $0.1$ and $1$, the two depletion controls,
a collisionless-baryon control, two particle samplings, finer time steps,
particle/smoothing refinement, and the shorter contact range. In the
collisionless-baryon control, the gas-labelled component starts with isotropic
velocity support; it is not a pressure switch applied to an identical fluid
initial state.

The comparison epoch was declared before the runs: first galaxy-core
pericentre plus two dynamical units. Core centres follow stellar elements
initially inside their cluster's scale radius. Maps project along $z$ and
apply the same Gaussian width $0.3a_1$. Tagged-component peaks and centres
of mass are diagnostic outputs. The observable source map includes all
particles without cohort labels. A relative gas-emission proxy weights each
gas element by $m_i\rho_i\sqrt{u_i}$; no instrument response, metallicity,
source-redshift distribution or measured lensing covariance is included.

For illustration, $M_1=10^{15}M_\odot$ and $a_1=300\,$kpc give
$v_*=3786\,$km\,s$^{-1}$ and $t_*=77.47\,$Myr. The initial relative speed
would be $3029\,$km\,s$^{-1}$, and the displayed post-pericentre interval
$154.9\,$Myr. Rates $\nu=0.1$ and $1$ correspond to $1.291$ and
$12.91\,$Gyr$^{-1}$. These units supply scale, not a fit to a named cluster.

\section{Conservation and formal checks}
The complete evolved Newtonian ledger is
\begin{equation}
 E_{\rm tot}=\sum_i\frac12m_iv_i^2+U_G+
 \sum_{i\in g}m_i u_i+\sum_{i\in C}(U_i+E_{B,i}).
\end{equation}
The largest error across the completed runs, normalized by initial
gravitational binding energy, was $6.84\times10^{-5}$, or $0.00684\%$.
The zero-contact error changed from $2.619\times10^{-5}$ to
$6.534\times10^{-6}$ when the step was halved, giving convergence order
$2.003$. Stochastic encounter histories do not permit the same deterministic
order inference for the nonzero-contact cases.

Maximum total momentum and angular-momentum discrepancies were below
$3.7\times10^{-15}$ and $4.3\times10^{-16}$ in simulation units.
An independent 10,000-pair calculation checked recoil, torque and the reduced-mass
work identity. The gas force and energy derivatives satisfy the instantaneous
power identity, and the sampled gas entropy
$\sum_i m_i\log[u_i/\rho_i^{\gamma_g-1}]$ increased in every hydrodynamic run.
All gas internal energies stayed positive; work and bath boundaries held.

Eight new Lean certificates check the reciprocal velocity update, its relative
speed and reduced-mass work, central torque cancellation, SPH pair power,
nonnegative viscous heat under its sign assumptions, register recharge for
negative work, and the absorbing zero-clock gate. Their axiom audit uses
standard logical axioms only. They supplement the earlier reciprocity and
finite-work identities. They do not certify the numerical PDE solution or
select the physical interaction. The standalone Lean/mathlib 4.19 runtime was
used; upstream 4.35 was not rebuilt.

\section{Encounter outcomes}
\begin{figure}[p]
\centering\includegraphics[width=\textwidth]{merger_maps.png}
\caption{Synthetic encounter maps at the declared epoch. Cyan contours show
total projected mass; white points show stellar tracers. Colour is the gas
emission proxy normalized to each panel's own maximum, so colour brightness
cannot compare absolute luminosities across panels. The first three panels
use 2688 elements; the depleted control uses 1344. These are model outputs,
not reconstructed observations.}
\end{figure}

The weak-contact encounters retain two separated total-mass concentrations.
The smaller cluster's stock peak remains close to its stellar peak, while
its tagged gas is displaced. The feature persists under the tested sampling
and resolution changes. In contrast, treating the gas-labelled component as
collisionless leaves it close to the stock.

\begin{table}[ht]
\centering\small
\begin{tabular}{@{}lrrrr@{}}
\toprule
Case & $t/t_*$ & $d_{C,g}/a_1$ & $d_{C,*}/a_1$ & Core separation/$a_1$\\
\midrule
No contact, baseline & 10.0 & 2.136 & 0.112 & 2.507\\
No contact, finer step & 10.0 & 2.136 & 0.112 & 2.507\\
No contact, more particles & 9.8 & 1.945 & 0.071 & 2.350\\
No contact, second sampling & 9.8 & 2.081 & 0.071 & 2.390\\
Weak contact, baseline & 10.0 & 2.214 & 0.180 & 2.394\\
Weak contact, more particles & 9.8 & 2.055 & 0.100 & 2.350\\
Weak contact, second sampling & 10.0 & 2.326 & 0.112 & 2.428\\
Strong contact, baseline & 10.8 & 3.086 & 0.250 & 1.425\\
Strong contact, finer step & 10.4 & 3.637 & 0.453 & 1.611\\
Strong contact, more particles & 10.0 & 2.907 & 0.206 & 1.626\\
Strong contact, shorter range & 10.8 & 2.780 & 0.212 & 1.379\\
Nearly depleted, $\nu=1$ & 10.0 & 2.088 & 0.071 & 2.499\\
Collisionless baryons & 10.0 & 0.071 & 0.050 & 2.577\\
\bottomrule
\end{tabular}
\caption{Distances between projected peaks of the smaller cluster's tagged
components. These are not measured lensing offsets. Each run uses its own
first-pericentre time plus the same fixed delay.}
\end{table}

For the weak-contact cases, $d_{C,g}/a_1=2.06$--$2.33$ and
$d_{C,*}/a_1=0.10$--$0.18$. With the illustrative length conversion these
are about $617$--$698\,$kpc and $30$--$54\,$kpc. They show the scale of
this chosen encounter; they should not be quoted as a prediction for the
Bullet Cluster. The two samplings and two particle resolutions support the
qualitative separation, while the map and offset magnitudes remain insufficiently
converged for an observational likelihood.

Strong contact also leaves separated concentrations in these examples, but
it changes core trajectories, gas morphology and the register energy much
more strongly. Its varying peak offsets cannot be read as a monotonic drag
constraint: halo structure and gas heating change along with momentum exchange.

\subsection{Depletion and recharge}
At the same $\nu=1$, the baseline active encounter accepts 7114 contacts;
the nearly depleted case accepts 61. Both preserve the work boundary.
The latter follows the zero-contact encounter closely at the comparison epoch.
With every register exactly at its floor, no contact is accepted and all
saved particle positions are identical to the zero-contact run.

Depletion is not a one-way loss of available work. Negative $W$ in
\eqref{eq:ledger} stores kinetic energy back in the register. The active
encounter's net work to mechanical motion is $-0.08490$; its registers gain
that energy. The nearly depleted case gains $0.003158$ and can reactivate
some initially slow elements. The exactly resting state remains absorbing
because the initial factor in~\eqref{eq:gate} vanishes.

These statements concern the specified internal-register implementation.
The calculation does not establish that real cluster stock begins depleted,
or that every component of the capacity medium has the same register state.

\section{The isolated controls constrain the target}
An encounter can retain two mass peaks even when the chosen contact does
not preserve an equilibrium halo. Isolated tests separate these questions.
At six dynamical units, the median radii relative to their initial values are
\begin{center}
\begin{tabular}{@{}lrrr@{}}
\toprule
Contact & Stock & Stellar tracers & Gas\\
\midrule
None & 0.968 & 0.957 & 0.921\\
Weak, $\nu=0.1$ & 0.920 & 0.969 & 1.024\\
Strong, $\nu=1$ & 0.727 & 1.269 & 1.464\\
\bottomrule
\end{tabular}
\end{center}
The zero-contact changes quantify finite sampling, truncation and gas/gravity
smoothing in the initial construction. Weak contact produces a modest additional
stock contraction over this interval. Strong contact contracts the stock by
about $27\%$ and expands the gas by about $46\%$, before any encounter is
introduced. Its target is therefore not an equilibrium prescription for the
declared Plummer halo.

\begin{figure}[ht]
\centering\includegraphics[width=\textwidth]{isolated_controls.png}
\caption{Isolated controls show the structural effect of the specified
contact. The changes at strong coupling prevent treating the initial profile
as an established stationary solution of that model.}
\end{figure}

This result sets a direct requirement on the next target calculation:
the selected kinetic--fluid interaction must preserve the intended isolated
equilibrium, or predict its controlled secular evolution, before its merger
offsets can constrain a universal rate. A small rate can keep changes modest
during one encounter; it does not by itself establish a galaxy RAR equilibrium
over an assembly history.

\section{Cohort accounting in a finite aperture}
Particle masses and parent labels are retained throughout the encounter.
For one parent cohort with initial baryonic mass $M_{b,L}$ and stock
$kM_{b,L}$, define retained fractions in an aperture by $b$ and $s$. Its
own material satisfies
\begin{equation}
 \frac{M_C}{M_b}=k\frac{s}{b}.
\end{equation}
The global cohort budget remains fixed while a present-aperture baryon
fraction changes through transport. Material from the other cluster is
recorded separately and then included in the observable total mass.

An exploratory aperture audit illustrates this in the refined weak-contact
run. Inside radius $a_1$ around the smaller stellar core, the parent retains
$b=0.0894$ of its baryons and $s=0.7826$ of its stock. The aperture also
contains $0.01091$ stock and $0.01325$ baryonic mass from the other cluster,
in units of $M_1$. Including both cohorts gives baryon fraction $0.0781$.
At radius $2a_1$, the total baryon fraction is $0.1707$ because the enlarged
aperture contains more gas and imported material.

These apertures were examined after the encounter calculation as accounting
diagnostics. They are not $R_{500}$ or $R_{200}$ measurements, and this disturbed
system supplies no fit to the relaxed X-COP sample. The calculation demonstrates
explicit transport behind an aperture/cohort distinction without creating
extra stock or applying a cap per present baryon.

\section{What the calculation establishes}
The conservative stock route admits a live three-dimensional encounter with
mass--gas separation, finite local exchange and explicit recoil. Weak contact
and the tested nearly depleted contact preserve that separation. The work gate
has its intended absorbing floor, while finite surplus permits energy return
and reactivation. Strong contact exposes the need to select an equilibrium
target and its physical rate jointly.

A predictive cluster application still requires the initial stock distribution
from an assembly model, a physically fixed contact target and rate, matching
of the work/bath registers to the capacity stress, and substantially higher
resolution with observational projection and likelihoods. The pair-relative
target and finite range introduced here are effective assumptions. There is
no theorem equating them to the manuscript's earlier phase-weighted cluster
map or to the single-carrier galactic contact. No high-redshift $a_0$ issue is
resolved by this calculation.

The research directory retains all executed sources, initial states,
histories, comparisons, formal audits and the interrupted first slow baseline.
The latter was replaced by an algebraically identical optimized force
implementation, whose pair forces were checked independently. One early
source-hash metadata entry was corrected from its saved start-of-run source;
the original record is retained. The preceding research and all four current
publication sources were preserved.

\begin{thebibliography}{9}
\bibitem{Clowe} D. Clowe et al., ``A direct empirical proof of the existence
of dark matter,'' Astrophys. J. 648, L109 (2006),
\href{https://arxiv.org/abs/astro-ph/0608407}{arXiv:astro-ph/0608407}.
\bibitem{Springel} V. Springel and G. R. Farrar, ``The speed of the `bullet'
in the merging galaxy cluster 1E0657-56,'' Mon. Not. R. Astron. Soc. 380,
911 (2007), \href{https://arxiv.org/abs/astro-ph/0703232}{arXiv:astro-ph/0703232}.
\bibitem{Cossins} P. J. Cossins, ``Smoothed Particle Hydrodynamics,''
\href{https://arxiv.org/abs/1007.1245}{arXiv:1007.1245} (2010).
\end{thebibliography}
\end{document}
